% Figure: the vapour-compression refrigeration cycle on a temperature-entropy
% diagram. A rounded dome stands for the two-phase region; the cycle runs
% 1->2 (compression, entropy rises slightly for a real compressor), 2->3
% (condensation at the high pressure, leftward along the high isotherm),
% 3->4 (throttling through the valve, an isenthalpic step that moves to the
% right in entropy at falling temperature), 4->1 (evaporation at the low
% pressure, rightward along the low isotherm). The rightward shift across the
% valve is the throttling irreversibility, the loss the energy balance misses.
\begin{figure}[htbp]
\centering
\begin{tikzpicture}
\begin{axis}[
  width=11cm, height=7.5cm,
  xlabel={specific entropy $s$ (arbitrary units)},
  ylabel={temperature $T$ (arbitrary units)},
  xmin=0, xmax=10, ymin=0, ymax=10,
  axis lines=left,
  xtick=\empty, ytick=\empty,
  tick label style={font=\scriptsize},
  label style={font=\small},
  clip=false,
]
% two-phase dome (a smooth bell as a stand-in saturation curve)
\addplot[enggray, thick, domain=1.5:8.5, samples=120]
  {2.0 + 6.0*exp(-((x-5)^2)/4.5)};
\node[enggray, font=\scriptsize, anchor=south] at (axis cs:5,8.2) {two-phase dome};
% high-pressure isotherm (condenser level) and low-pressure isotherm (evaporator)
\draw[enggray, dotted] (axis cs:0,6.5) -- (axis cs:10,6.5);
\draw[enggray, dotted] (axis cs:0,3.0) -- (axis cs:10,3.0);
\node[enggray, font=\scriptsize, anchor=west] at (axis cs:8.6,6.5) {$T_H$};
\node[enggray, font=\scriptsize, anchor=west] at (axis cs:8.6,3.0) {$T_C$};
% state points: 1 (evaporator exit, sat vapour, low T), 2 (compressor exit, high T),
% 3 (condenser exit, sat liquid, high P), 4 (valve exit, low T two-phase)
% 1 -> 2 compression (entropy rises a little for a real machine)
\draw[engblue, very thick, -{Stealth}] (axis cs:6.7,3.0) -- (axis cs:7.3,7.6);
% 2 -> 3 condensation, leftward along high isotherm down to subcooled liquid
\draw[engblue, very thick, -{Stealth}] (axis cs:7.3,7.6) -- (axis cs:2.6,6.5);
% 3 -> 4 throttling: isenthalpic, entropy increases (rightward), T falls
\draw[engred, very thick, -{Stealth}] (axis cs:2.6,6.5) -- (axis cs:4.6,3.0);
% 4 -> 1 evaporation along low isotherm rightward
\draw[engblue, very thick, -{Stealth}] (axis cs:4.6,3.0) -- (axis cs:6.7,3.0);
% labels
\node[font=\scriptsize, anchor=north] at (axis cs:6.7,2.9) {$1$};
\node[font=\scriptsize, anchor=south west] at (axis cs:7.3,7.6) {$2$};
\node[font=\scriptsize, anchor=south] at (axis cs:2.6,6.7) {$3$};
\node[font=\scriptsize, anchor=north] at (axis cs:4.6,2.9) {$4$};
\node[engred, font=\scriptsize, anchor=west, align=left] at (axis cs:3.2,4.8)
  {throttling:\\ entropy gain};
\node[engblue, font=\scriptsize, anchor=west] at (axis cs:7.0,5.3) {compressor};
\node[engblue, font=\scriptsize, anchor=south] at (axis cs:5.0,3.05) {evaporator};
\node[engblue, font=\scriptsize, anchor=south] at (axis cs:4.8,6.55) {condenser};
\end{axis}
\end{tikzpicture}
\caption[Refrigeration cycle on a temperature-entropy diagram]{The
vapour-compression refrigeration cycle drawn on temperature-entropy axes, with the
two-phase dome sketched as a smooth saturation curve. The refrigerant leaves the
evaporator as saturated vapour at state $1$, is compressed to the high pressure at
state $2$ (the slight rightward lean of the compression marks the real
compressor's entropy generation), condenses leftward along the high isotherm to
the saturated liquid at state $3$, and is throttled through the expansion valve to
state $4$. The throttle is adiabatic and isenthalpic, so an energy balance records
nothing across it, yet the step moves sharply to the right: the valve generates
entropy by dropping the pressure without recovering any work. That rightward shift
is the throttling irreversibility, the exergy the cycle destroys for the
convenience of a cheap valve in place of an expander. The cooling delivered is the
heat absorbed along the low isotherm $4\!\to\!1$; the work paid is the area the
real cycle fails to enclose against its reversible reference.}
\label{fig:vol04ch04:refrigeration-ts}
\end{figure}
